\section*{Chapter \ref{ch:iv:stochastic-calculus} --- Stochastic Calculus}

\begin{solution}{iq:iv:stochastic-calculus:1}
$d(W_t^2) = 2W_t\,dW_t + dt$. Integrating and taking expectations, the stochastic integral has mean zero, so $\E[W_t^2]
= t$. Without the $dt$ term the answer would be zero, which is impossible for a non-negative variable.
\iqlookfor{the second-order Itô term and the \omterm{def:iv:phone-and-technical-screens:sanity}{sanity check} that exposes its absence.}
\end{solution}

\begin{solution}{iq:iv:stochastic-calculus:2}
$d(W_t^3) = 3W_t^2\,dW_t + 3W_t\,dt$: the drift $3W_t$ is not zero, so it is not a martingale. $W_t^3 - 3\int_0^t
W_s\,ds$ is (its differential is $3W_t^2\,dW_t$, and the integrand is square-integrable).
\iqlookfor{reading the drift off Itô's formula and compensating it.}
\end{solution}

\begin{solution}{iq:iv:stochastic-calculus:3}
$(W_1, W_2)$ is Gaussian with correlation $\sqrt{1/2}$. For a centred bivariate normal, $\P(X > 0, Y > 0) = \tfrac14 +
\arcsin(\rho)/(2\pi) = \tfrac14 + \tfrac{\pi/4}{2\pi} = \tfrac38$. A simulation of 400\,000 pairs agrees.
\iqlookfor{the correlation of Brownian values and the orthant probability.}
\end{solution}

\begin{solution}{iq:iv:stochastic-calculus:4}
It is a linear functional of a Gaussian process, so Gaussian, with mean 0 and variance
\[
\int_0^T\!\!\int_0^T \min(s,u)\,ds\,du = \frac{T^3}{3}.
\] (Equivalently, $\int_0^T W_s\,ds = \int_0^T (T - s)\,dW_s$ by integration by parts, whose variance is $\int_0^T
(T - s)^2\,ds$.)
\iqlookfor{Gaussianity and the covariance integral, or the integration-by-parts representation.}
\end{solution}

\begin{solution}{iq:iv:stochastic-calculus:5}
By \cref{prop:iv:stochastic-calculus:exit}, $\P = 2/(1 + 2) = \tfrac23$, and $\E[\tau] = 1 \times 2 = 2$.
\iqlookfor{optional stopping on $W$ and on $W^2 - t$.}
\end{solution}

\begin{solution}{iq:iv:stochastic-calculus:6}
With $\mu = 0.5$, $\sigma = 1$, $a = b = 1$: $(1 - e^{1})/(e^{-1} - e^{1}) = 1/(1 + e^{-1}) \approx 0.731$. The
martingale is $e^{-2\mu X_t} = e^{-X_t}$.
\iqlookfor{the exponential martingale for a drifted walk.}
\end{solution}

\begin{solution}{iq:iv:stochastic-calculus:7}
$2\P(W_1 \ge 1) = 2(1 - \Phi(1)) \approx 0.317$. A discrete simulation checks the maximum only at the grid points and misses
excursions between them, so it underestimates (\cref{fig:iv:stochastic-calculus:reflect}). Fix it with a Brownian-bridge
correction: between two grid values $x, y$ below the barrier, the path crosses with probability $\exp(-2(a - x)(a - y)/\Delta
t)$; draw that event, and the estimate is unbiased on any grid.
\iqlookfor{the reflection principle, the direction of the discretisation bias and the bridge correction.}
\end{solution}

\begin{solution}{iq:iv:stochastic-calculus:8}
Mean $e^{0.08 \times 10} \approx 2.23$; median $e^{(0.08 - 0.08) \times 10} = 1.00$. The mean is carried by a minority of very
good paths; the typical investor ends where she started. Quote the median and a range, not the mean.
\iqlookfor{the $-\sigma^2/2$ correction and its practical meaning.}
\end{solution}

\begin{solution}{iq:iv:stochastic-calculus:9}
Take $\theta = (\mu - r)/\sigma = (0.08 - 0.03)/0.2 = 0.25$ and $d\mathbb Q/d\mathbb P = \exp(-0.25 W_T - 0.03125T)$. Under
$\mathbb Q$, $W_t + 0.25t$ is a Brownian motion, so $W$ has drift $-0.25$, and $dS = rS\,dt + \sigma S\,dW^{\mathbb Q}$.
\iqlookfor{the market price of risk as the Girsanov kernel.}
\end{solution}

\begin{solution}{iq:iv:stochastic-calculus:10}
$d_1 = \sigma\sqrt T/2 = 0.1$, $d_2 = -0.1$. Cash-or-nothing: $\Phi(-0.1) \approx 0.4602$. Asset-or-nothing: $100\,\Phi(0.1)
\approx 53.98$. Call $= 53.98 - 100 \times 0.4602 \approx 7.97$, the Black--Scholes value (and close to $0.4 \times 100
\times 0.2 = 8$).
\iqlookfor{digitals under two measures and their difference as the call.}
\end{solution}

\begin{solution}{iq:iv:stochastic-calculus:11}
By Feynman--Kac with $dX = dW$, $u(t, x) = \E[(x + W_T - W_t)^2] = x^2 + (T - t)$. Check: $u_t = -1$, $\tfrac12 u_{xx} = 1$,
sum 0; and $u(T, x) = x^2$.
\iqlookfor{writing the solution as an expectation and verifying it in the PDE.}
\end{solution}

\begin{solution}{iq:iv:stochastic-calculus:12}
$X_\tau = 1$ and $X_t - \mu t$ is a martingale; optional stopping (with $\E[\tau] < \infty$ for $\mu > 0$) gives $1 - \mu\E[\tau] =
0$, so $\E[\tau] = 1/\mu$. As $\mu \to 0$ the expectation tends to infinity: driftless Brownian motion reaches $+1$ with
probability one, but its hitting time has infinite mean (its density decays like $t^{-3/2}$). Certain and fast are different
things.
\iqlookfor{Wald for Brownian motion with drift and the heavy tail of the driftless hitting time.}
\end{solution}

\begin{solution}{iq:iv:stochastic-calculus:13}
$\tau$ is finite almost surely but $\E[\tau] = \infty$, and $W_{t \wedge \tau}$ is not uniformly integrable (it can go
arbitrarily negative before hitting 1), so the optional stopping theorem does not apply. With a lower barrier $-b$, the
stopped process is bounded, optional stopping holds, and $0 = \E[W_\tau] = p \times 1 - (1 - p)b$, so the chance of hitting 1
first is $b/(1 + b)$; it tends to one as $b \to \infty$ while the loss in the other case grows, which reconciles the two
statements.
\iqlookfor{the conditions of optional stopping and a bounded version that works.}
\end{solution}
